\documentclass[12pt, letterpaper]{article}

%-------Packages---------
\usepackage{newpxtext,newpxmath} % Palatino-like text & math (modern)
\usepackage{bboldx}
\usepackage{mathtools}
\usepackage{tikz}
\usetikzlibrary{calc,intersections}
\usepackage{tikz-cd}
\usepackage[all,arc]{xy}
\usepackage{graphicx}

\usepackage{enumerate}
\usepackage[dvipsnames]{xcolor}
\usepackage{float}
\usepackage[left=1in,top=1in,right=1in,bottom=1in]{geometry}
\usepackage[nocheck]{fancyhdr}
\usepackage{multicol}
\usepackage{setspace}

\usepackage{dsfont}
\usepackage{mathrsfs}

\usepackage{tcolorbox}
\tcbuselibrary{theorems,skins,breakable}

\usepackage{xparse}
\usepackage{import}
\usepackage[utf8]{inputenc}

\usepackage{hyperref}
\usepackage{cleveref}

\usepackage{xcolor, ifthen, tcolorbox}
\tcbuselibrary{theorems,skins,breakable}
% Theorem System
% The following environments are provided:
%   New Problem:        \begin{problem} ...
%   New Question Header:   \qh (then followed by subproblems)
%   Subproblem:     \begin{subproblem} ...
%   Problem Proof:  \begin{problemproof} ...
%   Proof:          \\begin{pf}{color} ...
%   Remark:         \rmk (sentence), \rmkb (block)

% Define custom colors
\definecolor{thmscol}{HTML}{F4565B} %For Problems
\definecolor{lemscol}{HTML}{FE844C} %For Lemmes
\definecolor{prpscol}{HTML}{00A7EF} %For proposition
\definecolor{corscol}{HTML}{23DAFF} %For corrolaries
\definecolor{rmkscol}{HTML}{6DEEC5} %For remarks
\definecolor{defscol}{HTML}{18C991} %For definitions
\definecolor{exmscol}{HTML}{7DB800} %For examples
\definecolor{clmscol}{HTML}{165c58} %For claims

%Change between dark(black)/light(white) mode
\newcommand{\colormode}{white}
\ifthenelse{\equal{\colormode}{white}}
      {\newcommand{\TextColor}{black}}
      {\newcommand{\TextColor}{white}}
\newcommand{\pbcolormain}{thmscol!80!\colormode}
\newcommand{\pbcolorsecond}{thmscol!38!\colormode}


% ============================
% Problem Counter
% ============================
\newcounter{subproblemcounter}
\newcounter{problemcounter}

\newcommand{\q}{
    \setcounter{subproblemcounter}{0}%
    \stepcounter{problemcounter} % Increment problem counter
    \color{\TextColor}Problem \arabic{problemcounter}: % Display the problem counter
}

\newcommand{\stepsub}{
    \stepcounter{subproblemcounter}%
    \color{\TextColor}(\alph{subproblemcounter})
}
% ============================
% Problem
% ============================
\newtcolorbox{pb}[1]{
  enhanced,
  frame hidden,
  titlerule=0mm,
  toptitle=1mm,
  bottomtitle=1mm,
  fonttitle=\bfseries\large,
  coltitle=black,
  colbacktitle=\pbcolormain,
  colback=\pbcolorsecond,
  title={\q #1},
  coltext=\TextColor
}

\newtcolorbox{spb}[1]{
  enhanced,
  frame hidden,
  titlerule=0mm,
  toptitle=1mm,
  bottomtitle=1mm,
  fonttitle=\bfseries\large,
  coltitle=black,
  colbacktitle=\pbcolormain,
  colback=\pbcolorsecond,
  title={\stepsub #1},
  coltext=\TextColor,
  attach boxed title to top left={xshift=3mm,yshift=-2mm},
  boxed title style={
    colback=\pbcolormain,
    colframe=\pbcolormain,
  }
}

\newtcolorbox{pbh}[1]{
    enhanced,
    frame hidden,
    titlerule=0mm,
    toptitle=1mm,
    bottomtitle=1mm,
    fonttitle=\bfseries\large,
    coltitle=black,
    colbacktitle=\pbcolormain,
    colback=\pbcolorsecond,
    title={\q #1},
    boxrule=0pt,
    interior hidden,
    attach boxed title to top left={xshift=3mm,yshift=-2mm},
    boxed title style={
        colback=\pbcolormain,
        colframe=\pbcolormain,
    }
}

\newenvironment{problem}[1]{%
  \newpage
  \begin{pb}{#1}
    }{
  \end{pb}
}

\newenvironment{subproblem}[1]{%
  \begin{spb}{#1}
    }{
  \end{spb}
}

\newenvironment{problemproof}{%
    \begin{pf}{\pbcolormain}
        }{
    \end{pf}
}

\newcommand{\qh}[1][]{
    \newpage
    \begin{pbh}{#1}\end{pbh} % Display the problem counter
}

% ============================
% Claim
% ============================

\newtcbtheorem[number within=problemcounter]{claim}{Claim}
{
    enhanced,
    frame hidden,
    titlerule=0mm,
    toptitle=1mm,
    bottomtitle=1mm,
    fonttitle=\bfseries\large,
    coltitle=black,
    colbacktitle=clmscol!50!\colormode,
    colback=clmscol!20!\colormode,
}{clm}

\NewDocumentCommand{\clm}{m+m}{
    \begin{claim*}{#1}{}
        #2
    \end{claim*}
}

\newenvironment{clmpf}{
	{\noindent{\it \textbf{Proof.}}
	\setlength{\parindent}{0pt}
	}
	\tcolorbox[blanker,breakable,left=5mm,parbox=false,
    before upper={\parindent15pt},
    after skip=10pt,
	borderline west={1mm}{0pt}{clmscol!50!\colormode}]
}{
    \textcolor{clmscol!50!\colormode}{\hbox{}\nobreak\hfill$\blacksquare$}
    \endtcolorbox
}

\NewDocumentCommand{\clmp}{m+m+m}{
    \begin{claim*}{#1}{}
        #2
    \end{claim*}

    \begin{clmpf}
        #3
    \end{clmpf}
}


% ============================
% Proof
% ============================

\newenvironment{pf}[1]{%
  \def\pfcolor{#1}% Store the color in a macro
  \noindent{\it \textbf{Proof.}}%
  \tcolorbox[blanker,breakable,left=5mm,parbox=false,%
    before upper={\parindent15pt},%
    after skip=10pt,%
    borderline west={1mm}{0pt}{\pfcolor},%
    coltext=\TextColor
  ]%
  \setlength{\parindent}{0pt}%
}{%
  \textcolor{\pfcolor}{\hbox{}\nobreak\hfill$\blacksquare$}%
  \endtcolorbox%
}

\newtcolorbox{solution}{
  blanker,
  breakable,
  left=5mm,
  parbox=false,%
  before upper={\parindent15pt},%
  after skip=10pt,%
  borderline west={1mm}{0pt}{\pbcolormain},%
  coltext=\TextColor
}


% ============================
% Remark
% ============================


\NewDocumentCommand{\rmk}{+m}{
    {\it \color{rmkscol!100!white}#1}
}

\newenvironment{remark}{
    \par
    \vspace{5pt}
    \begin{minipage}{\textwidth}
        {\par\noindent{\textbf{Remark.}}}
        \tcolorbox[blanker,breakable,left=5mm,
        before skip=10pt,after skip=10pt,
        borderline west={1mm}{0pt}{rmkscol!80!\colormode},
        coltext=\TextColor
        ]
}{
        \endtcolorbox
    \end{minipage}
    \vspace{5pt}
}

\NewDocumentCommand{\rmkb}{+m}{
    \begin{remark}
        #1
    \end{remark}
}


% Define a new command for problems
\renewcommand{\headrule}{\color{\TextColor}\hrulefill}
\newcommand{\C}{\mathbb{C}}
\newcommand{\R}{\mathbb{R}}
\newcommand{\Q}{\mathbb{Q}}
\newcommand{\Z}{\mathbb{Z}}
\newcommand{\N}{\mathbb{N}}
\newcommand{\p}{\mathfrak{p}}
\newcommand{\F}{\mathbb{F}}
\newcommand{\contr}{\Rightarrow \Leftarrow}
\newcommand{\s}{\(\,\)}
\renewcommand{\t}[1]{\text{#1}}
\newcommand{\Spec}{\mathrm{Spec}}
\newcommand{\mf}[1]{\mathfrak{#1}}

% Meta data
\setstretch{1.2}

\begin{document}
\color{\TextColor}
\pagecolor{\colormode}
\setlength{\parindent}{0pt}
\pagestyle{fancy}
\fancyhf{}
\fancyhead[LOE]{\color{\TextColor}\slshape{\textbf{Vincent Ferrara}}}
\fancyhead[ROE]{\color{\TextColor}HW3 --- \thepage}
\fancyhead[C]{\color{\TextColor}Math 614}

\begin{problem}{}
    Let $R$ be a ring, $S \subset R$ a multiplicative system, and $M$ an $R$-module.
\end{problem}
\begin{subproblem}{}
    Formulate and prove a universal property for the natural $R$-module homomorphism $M \to S^{-1}M$, analogous to our universal
    property for the localization homomorphism $R \to S^{-1}R$ for rings.

    (Universal Property):

    Let $N$ be a $R$-module s.t. every $s \in S$ acts invertibly (i.e. $n \mapsto sn$ is an automorphism).

    Then for every $R$-linear map $f: M \to N$ there exists a unique $S^{-1}R$-linear map $g: S^{-1}M \to N$ s.t. the diagram commutes
    \[\begin{tikzcd}
	M & {S^{-1}M} \\
	& N
	\arrow["\pi", from=1-1, to=1-2]
	\arrow["f"', from=1-1, to=2-2]
	\arrow["g", dashed, from=1-2, to=2-2]
    \end{tikzcd}\]
\end{subproblem}
\clmp{}{
    For $S^{-1}R$ a localization of $R$ if $N$ is a $R$-module where $s \in S$ acts invertibly on $N$ then the action
    \[(r/s)n = h_s(rn) = n(r/s)\]
    where $h_s$ is the inverse map of $n \to sn$ makes $N$ a $S^{-1}R$ module.
}{
    \setlength\parindent{0pt}
    Well-Defined:

    If $r/s=r'/s'$ then there exists a $t \in S$ s.t. $t(rs'-r's)=0$. So for $n$ in $N$
    \[0=t(rs'-r's)n=t(rs'n-r'sn)\]

    Thus since $t$ acts invertibly,
    \begin{align*}
        rs'n&=r'sn\\
        rn&=sh_{s'}(r'n)\\
        h_{s}(rn)&=h_{s'}(r'n)
    \end{align*}

    Thus, $h_{s}(rn)$ only depends on the equivalence class of $r/s$, so the action is well-defined.

    Module:

    Let $r/s, r'/s' \in S^{-1}R$ and $n,m \in N$ then
    \begin{enumerate}
        \item \[(r/s)(n+m)=h_{s}(r(n+m))=h_{s}(rn+rm)=h_{s}(rn)+h_{s}(rm)=(r/s)n+(r/s)m\] 
        \item \[(r/s+r'/s')n=(\frac{rs'+r's}{s's})n=h_{s's}(rs'n+r'sn)\]
    \[=h_{s}(h_{s'}(rs'n+r'sn))=h_s(rn)+h_{s'}{r'n}=(r/s)n+(r'/s')n\]
        \item \[(r/s \cdot r'/s')n=(rr'/ss')n=h_{s'}(r'h_{s}(rn))=h_{s'}(r'(r/s)n)=(r'/s')((r/s)n)\]
        \item \[1/1 \cdot n= h_{1}(n)=n\]
    \end{enumerate}
    Therefore, $N$ can be seen as a $S^{-1}R$ module under this action, and so any $R$-module which has all elements $s \in S$ act invertibly on it this implies it is a $S^{-1}R$-module.
}
\begin{problemproof}
    Existence: Shown in class $S^{-1}M$ is a well-defined $S^{-1}R$ module where all $s \in S$ act invertibly.

    This is true since for $m/t \in S^{-1}M$ we have that $s^{-1}s(m/t)=m/t$ and $ss^{-1}(m/t)=mt$

    Let $N$ be a $R$-modules s.t. every $s \in S$ acts invertibly.

    Let $f : M \to N$ be a $R$-linear map.

    Define: $g : S^{-1}M \to N$ s.t. $g(m/s)=h_s^{-1}f(m)=f(m)/s$ under the action defined above where $h_s^{-1}$ is the inverse of the map $n \mapsto sn$.

    Well-Defined: Suppose $m/s=m'/s'$ in $S^{-1}M$. Then there exists a $t \in S$ s.t.
    \[t(s'm-sm')=0\]
    Now apply $f$\[\begin{tikzcd}
	M & {S^{-1}M} \\
	& N
	\arrow["\pi", from=1-1, to=1-2]
	\arrow["f"', from=1-1, to=2-2]
	\arrow["g", dashed, from=1-2, to=2-2]
\end{tikzcd}\]
    \[t(s'f(m)-sf(m'))=0\]
    Since $t$ acts invertible apply its inverse, and we get
    \begin{align*}
        s'f(m)&=sf(m')\\
        h_s(s'f(m))&=h_s(sf(m'))\\
        s'h_s(f(m))&=f(m')\\
        h_{s'}s'(h_s(f(m)))&=h_{s'}(f(m'))\\
        h_s(f(m))&=h_{s'}(f(m'))\\
        g(m/s)&=g(m'/s')
    \end{align*}

    Homomorphism:

    Let $m/s,m'/s' \in S^{-1}M$

    Since $s'sn=s'(sn) \implies h_{s's}=h_{s}\circ h_{s'}$

    \[g(\frac{m}{s}+\frac{m'}{s'})=g(\frac{s'm+sm'}{ss'})=h_{ss'}(s'f(m)+sf(m'))\]
    \[=h_{s'}(s'h_{s}(f(m)))+h_{s'}(h_{s}(f(m')))=h_{s}(f(m))+h_{s'}(f(m'))=g(m/s)+g(m'/s')\]

    Let $m/s \in S^{-1}M$ and $r/t \in S^{-1}R$.
    \begin{align*}
        g(r/t \cdot m/s)=g(rm/ts)=g(rm/st)=h_{st}(f(rm))=h_{t}(h_{s}(rf(m)))\\
        =h_{t}(rh_s(f(m)))=(r/t)(g(m/s))
    \end{align*}

    Thus, $g$ is $S^{-1}R$-linear.

    Diagram:
    
    Let $m \in M$ then $g(\pi(m))=g(m/1)=h_{1}(f(m))=f(m)$. Thus, $g \circ \pi = f$.

    Uniqueness:

    Let $h : S^{-1}M \to N$ be $S^{-1}R$-linear with $h \circ \pi = f$. For $m/s \in S^{-1}M$,
    \[h(m/s)=(1/s)h(\pi(m))=(1/s)f(m)=h_{s}f(m)=g(m/s)\]
    Thus, $g$ is unique.
\end{problemproof}
\begin{subproblem}{}
    Prove that if $M$ is an $R$-module such that every element $u \in S$ acts invertibly on
$M$, then there is an isomorphism $M \cong S^{-1}M$ of $R$-modules.
\end{subproblem}
\begin{problemproof}
    Assume $M$ is an $R$ modules s.t. for all $s \in S$ $s$ acts invertibly on $M$. Let $h_s$ be the inverse of $m \mapsto sm$.

    Define $\pi : M \to S^{-1}M$ s.t. $\varphi(m)=m/1$

    By the claim you can view $M$ as a $S^{-1}R$ module with the action $(r/s)m=h_s(rm)$.

    Using the universal property this implies that there is a unique map $S^{-1}R$-linear map s.t. $\psi : S^{-1}M \to M$ s.t. $\psi \circ \pi = \mathrm{id}$.
    
    Again using the universal property with $S^{-1}M$ we get there exists a unique map $\theta : S^{-1}M \to S^{-1}M$ s.t. $\theta \circ \pi = \pi$

    Since $\mathrm{id}$ also works this implies that $\theta = \mathrm{id}$

    Also, $(\pi \circ \varphi) \circ \pi = \pi$ also works this implies $\pi \circ \varphi =\mathrm{id}$.

    Thus, $\pi$ and $\varphi$ are $S^{-1}R$-linear maps and are inverses.

    Thus, $S^{-1}M \cong M$ as $S^{-1}R$ modules.

    But the same maps work for $R$ modules since $\varphi(rm/s)=r\varphi(m/s)$ and $\pi(rm)=rm/1=(r)(m/1)$.

    Thus, $S^{-1}M \cong M$ as $R$ modules.
\end{problemproof}
\begin{subproblem}{}
    Let $f : M \to N$ be an $R$-module homomorphism. Show that there is a natural 
    $S^{-1}R$-module homomorphism $S^{-1}f : S^{-1}M \to S^{-1}N$ induced by $f$, given by 
    $\frac{m}{s} \to \frac{f(m)}{s}$ for $s \in S$ and $m \in M$.
\end{subproblem}
\begin{problemproof}
    Well-Defined:

    Let $m/s=m'/s'$ this implies there exists a $t\in S$ s.t. $t(s'm+sm')=0$

    $0=f(t(s'm+sm'))=t(s'f(m)+sf(m'))$ since $f$  is $R$-linear.

    Therefore, this implies $f(m')/s'=f(m)/s$ thus well-defined.

    Homomorphism:

    Let $m/s, m'/s' \in S^{-1}M$

    $S^{-1}f(m/s+m'/s')=S^{-1}f(\frac{ms'+m's}{ss'})=f(ms'+m's)/s's=s'f(m)/s's+sf(m')/s's=f(m)/s+f(m')/s=S^{-1}f(m/s)+S^{-1}f(m'/s')$.

    Let $r/t \in S^{-1}R$

    $S^{-1}f(r/t \cdot m/s)=S^{-1}f(rm/ts)=rf(m)/ts=(r/t) \cdot f(m)/s = (r/t) \cdot S^{-1}f(m/s)$

    Thus, $S^{-1}f$ is a $S^{-1}R$-module homomorphism.
\end{problemproof}

\begin{problem}{}
    
\end{problem}
\stepcounter{subproblemcounter}
\begin{subproblem}{Localization is Exact}
    Let $S \subset R$ be a multiplicative system. Given an exact sequence
    $L \overset{f}{\longrightarrow} M \overset{g}{\longrightarrow} N$, prove the sequence of $S^{-1}R$-module homomorphisms
    \[S^{-1}L \overset{S^{-1}f}{\longrightarrow} S^{-1}M \overset{S^{-1}g}{\longrightarrow} S^{-1}N\]
    is also exact.
\end{subproblem}
\begin{problemproof}
    Assume $L \overset{f}{\longrightarrow} M \overset{g}{\longrightarrow} N$ is an exact sequence of modules.

    Let $x \in \mathrm{Im}(S^{-1}f)$ thus $f(l)/s=x$ for some $l \in L$.

    So, since the sequence is exact this implies that $f(l) \in \ker(g)$. Thus, $S^{-1}g(f(l)/s)=0/s=0 \implies x \in \ker(S^{-1}g)$.
    
    Let  $x \in \ker(S^{-1}g)$ then $x=m/s$ for $m \in M$ and $s \in S$, so $S^{-1}g(m/s)=g(m)/s=0$. This implies there exists a $t \in S$ s.t. $tg(m)=0$. Since $g$ is
    a $R$-module homomorphism this implies that $tg(m)=g(tm)$. Thus, $tm \in \ker{g} = \mathrm{Im}(f)$. So there is an $l \in L$ s.t. $f(l)=tm$. So,
    \[\frac{m}{s}=\frac{tm}{ts}=\frac{f(l)}{ts}=S^{-1}f\left( \frac{l}{ts} \right)\]
    This shows $m/s \in \mathrm{Im}(S^{-1}f)$.

    Therefore, $\mathrm{Im}(S^{-1}f)=\ker(S^{-1}g)$, so this is an exact sequence.
\end{problemproof}

\begin{subproblem}{Exactness is a local property}
    For any sequence $L \overset{f}{\longrightarrow} M \overset{g}{\longrightarrow} N$ of $R$-modules homomorphisms,
    prove that the following are equivalent:
    \begin{enumerate}[i.]
        \item The sequence $L \overset{f}{\longrightarrow} M \overset{g}{\longrightarrow} N$ is exact.
        \item The sequence $L_\p \overset{f_\p}{\longrightarrow} M_\p \overset{g_\p}{\longrightarrow} N_\p$ is exact for all $\p \in \Spec(R)$.
        \item The sequence $L_\mf{m} \overset{f_\mf{m}}{\longrightarrow} M_\mf{m} \overset{g_\mf{m}}{\longrightarrow} N_\mf{m}$ is exact for all maximal ideals $\mf{m} \in \mathrm{mSpec}(R)$.
    \end{enumerate}
\end{subproblem}
\begin{problemproof}
    $(i.) \implies (ii.)$ from part (b). $(ii.) \implies (iii.)$ since maximal ideals are a subset of $\Spec(R)$.

    $(iii.) \implies (i.)$:

    Assume $(iii.)$

    Let $x \in \mathrm{Im}(f)$ then there exists an $l \in L$ s.t. $f(l)=x$. Let $A=\{r \in R \mid rg(x)=0\}$.

    This is an ideal since if $a,b \in A$ then $(a+b)g(x)=ag(x)+bg(x)=0$ and for $a \in A$ and $r \in R$ we have $ra g(x)=r\cdot 0 =0$, so $A$ is an ideal.

    Assume $A$ is a proper ideal then $A$ is in some maximal ideal $\mf{m}$. By assumption, we know that $g_m(f_m(l/1))=0$. This implies that $g(f(l))/1=0/1$, so there exists some $a \in R \setminus \mf{m}$ s.t. $ag(f(l))=0$ but this is a contradiction since $A$ is a subset of $\mf{m}$.

    Thus, $A$ is not proper and $A=R$, thus $g(x)=0 \implies x \in \ker{g}$.

    Let $x \in \ker{g}$. Let $A=\{r \in R \mid rx \in \mathrm{Im}(f)\}$. This is an ideal for similar reasons as above. Assume that this is a proper ideal.

    Thus, there exists a maximal ideal $\mf{m}$ that contains $A$.

    Since $x/1 \in \ker(g_\mf{m})=\mathrm{Im}f_\mf{m}$ by assumption, 
    there exists a $l / s \in L_\mf{m}$ s.t. $f_m(l/s)=f(l)/s=x/1$. 
    This implies there exists a $a \in R \setminus \mf{m}$ s.t.
    
    $a(f(l)-ns)=0 \implies f(al)=asx$, so $asn \in \mathrm{Im}(f)$. 
    Thus, $as \in A$, but since $a \in R \setminus \mf{m}$ and $s \in R \setminus \mf{m}$ this implies that $as \in R \setminus \mf{m}$ since $\mf{m}$ is prime, but this is a contradiction. Thus, $A=R$, so $1*x=x \in \mathrm{Im}(f)$.

    Thus, $\ker(g)=\mathrm{Im}(f)$, so the sequence is exact.
\end{problemproof}

\stepcounter{problemcounter}

\begin{problem}{}
    Let $(R, m)$ be a local domain (i.e. a local ring which is also a domain) which is not a field.
Let $K$ be the fraction field of $R$, which is naturally an $R$-algebra via the inclusion $R \subset K$.
Prove that $K$ is not a finite R-algebra.
\end{problem}
\begin{problemproof}
    Assume $K$ is a finite $R$-algebra. Proven in class this implies that $K$ is integral.

    This implies that for all $1/r \in K$ s.t. $r \neq 0\in R$ then there exists $r_0,\dots,r_{n-1} \in R$ s.t. $(1/r)^n+r_{n-1}(1/r)^{n-1}+\dots+r_0=0$.

    Multiplying by $r^n$ gives $1+r_{n-1}r+\dots+r_0r^n=0$. This implies that $1 \in (r) \implies (r)=R$. Thus, $r$ is a unit, but this implies that $R$ is a field since $r$ was arbitrary.

    This is a contradiction thus $K$ is not a finite $R$-algebra.
\end{problemproof}

\begin{problem}{}
    Let $R$ be a domain with fraction field K. The normalization $\widetilde{R}$ of $R$ is the integral closure
of $R$ in $K$, i.e. $\widetilde{R} = \{y \in K \mid y \t{ is integral over } R\}$, and we say that R is normal if $R = \widetilde{R}$.
Prove that any UFD is normal.
\end{problem}
\begin{problemproof}
    Let $x = a/b \in \widetilde{R}$ with $a,b \in R$ s.t. $b \neq 0$, and $\gcd(a,b)=1$.

    Then there exists $r_{n-1},\dots,r_0 \in R$ s.t. $x^n+r_{n-1}x^{n-1}+\dots+r_0=0$.
    
    Multiply by $b^n$
    \[a^n+r_{n-1}ba^{n-1}+\dots+r_0b^n=0\]
    \[a^n=-b(r_{n-1}a^{n-1}+\dots+r_0b^{n-1})\]
    This implies that $b \mid a^n$ in $R$. Since $R$ is a UFD factor $a,b$ into irreducibles. Then if $p \mid b$ for some irreducible $p$. Then $p \mid a^n$, so $p \mid a$ since $p$ is irreducible. But this is a contradiction since $\gcd(a,b)=1$, so $b$ is a unit and $x \in R$.

    Therefore, $\widetilde{R} = R$.
\end{problemproof}

\qh
\begin{subproblem}{}
    Show that $\C[x,y] \diagup (y^3-x^5)$ is a domain which is not normal, and compute its normalization.
\end{subproblem}
\begin{problemproof}
    $\C[x,y] \diagup (y^3-x^5) \cong \C[t^3,t^5]$.

    Consider $x^3-t^3=0$ for $x=t$. This implies $t \in \widetilde{\C[t^3,t^5]}$. Thus, $\C[t] \subset \widetilde{\C[t^3,t^5]}$, but since $\C[t]$ is a UFD this implies it is normal thus $\widetilde{\C[t^3,t^5]}=\C[t]$. So $\C[t^3,t^5]$ is not normal, with normalization $\C[t]$.
\end{problemproof}

\begin{subproblem}{}
    Compute the normalization of $\Z\left[ \sqrt{-3} \right]$
\end{subproblem}
\clmp{}{
    If $\alpha$ is integral over $R$ then $R[\alpha]$ is finite.
}{
    Assume $\alpha$ is integral over $R$. This implies there exists a $r_{n-1},\dots,r_0 \in R$ s.t.
    \[\alpha^n+r_{n-1}\alpha^{n-1}+\dots+r_0=0\]
    This implies that for $\alpha^k$ for $k \geq n$ are a $R$-linear combination of $1,\dots,\alpha^{n-1}$, so $\{1,\dots,\alpha^{n-1}\}$ generate $R[\alpha]$
    as an $R$-module. Thus, $R[\alpha]$ is finite over $R$.
}
\begin{problemproof}
    Consider $\omega$ where $\omega$ is the cube root of unity.

    Since $\omega^2+\omega+1=0$ this implies that $\omega$ is integral over $\Z$, so $\Z[\omega]$ is integral over $\Z$ by the proof above.

    Since $\Z \subset \Z[\sqrt{-3}] \subset \Z[\omega]$ this implies $\Z[\omega]$ is integral over $\Z[\sqrt{-3}]$ since the same monic equation in $\Z$ will work in $\Z[\sqrt{-3}]$.

    Thus, $\Z[\omega] \subseteq \widetilde{\Z[\sqrt{-3}]}$
    
    Since $\Z[\omega]$ is a UFD by problem 5, this implies that $\Z[\omega]$ is normal thus integrally closed. Therefore, $\Z[\omega] = \widetilde{\Z[\omega]}$.

    Let $a \in \widetilde{\Z[\sqrt{-3}]}$. This implies that there exists $r_{n-1},\dots,r_0 \in \Z[\sqrt{-3}]$ s.t. $a^n+r_{n-1}a^{n-1}+\dots+r_0=0$. Since $\Z[\sqrt{-3}] \subset \Z[\omega] \implies r_i \in \Z[\omega]$.

    Thus, $a \in \widetilde{\Z[\omega]}=\Z[\omega]$.
    
    Therefore, $\Z[\omega]=\widetilde{\Z[\sqrt{-3}]}$.
\end{problemproof}

\qh[]
\begin{subproblem}{Normalization commutes with localization}
    Let $R$ be a domain and $S \subset R$ a multiplicative system.
    Via the inclusion $R \hookrightarrow \widetilde{R}$, we can also regard $S$ as a
    multiplicative system in the normalization $\widetilde{R}$. Prove that there is an isomorphism
    of rings
    \[S^{-1}\widetilde{R} \cong \widetilde{S^{-1}R},\]
    where $\widetilde{S^{-1}R}$ is the normalization of $S^{-1}R$.
\end{subproblem}
\begin{problemproof}
    Let $K = \mathrm{Frac}(S^{-1}R) \cong \mathrm{Frac}(R)$ since $R$ is a domain, and we are give $S^{-1}R$ is a domain. 
    
    Consider $\varphi : S^{-1}\widetilde{R} \hookrightarrow K$, so $S^{-1}\widetilde{R}$ is isomorphic to a subring of $K$ and $\widetilde{S^{-1}R} \subset K$.

    Let $A$ be the isomorphic subring.

    Since $\widetilde{R}$ is integral over $R$ this implies that for $x=r/s \in S^{-1}\widetilde{R}$. Thus $r \in \widetilde{R}$, so there exists a monic polnomial s.t. 
    $r^n+r_{n-1}r^{n-1}+\dots+r_0=0$ for $r_i \in R$.

    Thus, in $S^{-1}\widetilde{R}$, $(r^n/1)+r_{n-1}(r^{n-1}/1)+\dots+r_0=0 \implies x^n+r_{n-1}/s \cdot x^{n-1} + \dots + r_0/s^n=0$

    Thus, $S^{-1}\widetilde{R}$ is integral over $S^{-1}R$. This implies since $A$ is isomorphic that $A$ is integral over $S^{-1}R$ thus $A \subset \widetilde{S^{-1}R}$.

    Let $x \in \widetilde{S^{-1}R}$. This means that $x$ is integral over $S^{-1}R$, then
    
    $x^n+a_{n-1}x^{n-1}+\dots+a_0 = 0$ for $a_i \in S^{-1}R$.

    Let $a_i=c_i/s_i$. Let $d=\prod_{i=0}^{n-1}s_i \in S$ be a common denominator then multiply by $d^n$. We get

    $(dx)^n+c_{n-1}(dx)^{n-1} +  \dots + c_0=0$ for $c_i \in R$.

    Thus, $dx$ is integral over $R$, so $dx \in \widetilde{R} \implies sx/1 \cdot 1/s = x \in A$. Thus, $A=\widetilde{S^{-1}R}$.

    Therefore, $S^{-1}\widetilde{R} \cong A = \widetilde{S^{-1}R}$
\end{problemproof}

\begin{subproblem}{Normality is a local property}
    Let $R$ be a domain. Prove that the following are
equivalent:
\begin{enumerate}[i.]
    \item $R$ is normal
    \item $R_\p$ is normal for all $\p \in \Spec(R)$.
    \item $R_\mf{m}$ is normal all $\mf{m} \in \mathrm{mSpec}(R)$
\end{enumerate}
\end{subproblem}
\begin{problemproof}
    $(i) \implies (ii)$ by (a) since if $R$ is normal then $S^{-1}R = S^{-1}\widetilde{R} \cong \widetilde{S^{-1}R}$.

    $(ii) \implies (iii)$ since maximal ideals are a subset of prime ideals.

    $(iii) \implies (i)$:

    Assume $(iii)$.

    Let $x \in \widetilde{R}$ then there exists $r_0,\dots,r_{n-1} \in R$ s.t. $x^n+r_{n-1}x^{n-1}+\dots+r_0=0$.

    Let $I=\{a \in R \mid ax \in R\}$ where $ax$ is multiplication in $R$'s fraction field.

    Assume $I$ is proper, so there exists a maximal ideal $\mf{m}$ that contains $I$.

    Consider $R_{\mf{m}}$ we can view the relation $x$ has in $R_\mf{m}$ thus

    $x^n+r_{n-1}/1 \cdot x^{n-1}+\dots+r_0/1=0$ thus $x \in R_\mf{m}$ since by assumption $R_\mf{m}$ is normal.
    
    Let $x=r/s$ for $r \in R$ and $s \in R \setminus \mf{m}$.

    This implies that $(sx)^n+r_{n-1}s(sx)^{n-1}+\dots+s^nr_0=0$ $sx \in \widetilde{R}$ and $sx \in R$.

    Thus, $s \in I$ but this is a contradiction since $s \in R \setminus \mf{m}$ and $I$ is contained in $\mf{m}$. Thus, $I=R$ and $1x=x \in R$.
    
    Thus, $\widetilde{R}=R$.
\end{problemproof}
\end{document}